\documentclass[a4paper,landscape,headrule,footrule]{foils}

%\usepackage{nttfoilhead}
%\newcommand{\myslide}[1]{\foilhead[-25mm]{\raisebox{12mm}[0mm]{\emp{#1}}}}
%\newcommand{\myslider}[1]{\rotatefoilhead[-25mm]{\raisebox{12mm}[0mm]{\emp{#1}}}}
%\newcommand{\myslider}[1]{\rotatefoilhead{\raisebox{-8mm}{\emp{#1}}}}

%%
%%% macros for Theories of Grammar
%%%
\usepackage{polyglossia}
\setdefaultlanguage{english}
%\setmainfont{TeX Gyre Pagella}


\newcommand{\logo}{~}
\newcommand{\header}[3]{%
\title{\vspace*{-2ex} \large JMORF/MORS --- Morpho-Syntax
\\[2ex] \Large  \red{#2} \\ \red{#3}}
\author{\blu{Francis Bond}   \\ 
\normalsize  \textbf{Palacký University}\\
\normalsize  \url{https://fcbond.github.io/}\\
\normalsize  \texttt{bond@ieee.org}}
\MyLogo{JMORF (2026)}
\renewcommand{\logo}{#2}
\hypersetup{
   pdfinfo={
     Author={Francis Bond},
     Title={#1: #2},
     Subject={JMORF/MORS: Morpho-Syntax},
     Keywords={Syntax, Morphology, HPSG, Unification, Constructions},
     License={CC BY 4.0}
   }
 }

\date{#1
  \\ Location: KB-1.11}
}
\usepackage{hyperref}
\hypersetup{
     colorlinks,
     linkcolor={blue!50!black},
     citecolor={red!50!black},
     urlcolor={blue!80!black}
   }



\usepackage{xcolor}
\usepackage{graphicx}
\newcommand{\blu}[1]{\textcolor{blue}{#1}}
\newcommand{\red}[1]{\textcolor{red}{#1}}
\newcommand{\grn}[1]{\textcolor{green}{#1}}
\newcommand{\hide}[1]{\textcolor{white}{#1}}
\newcommand{\emp}[1]{\textcolor{blue}{#1}}
\newcommand{\txx}[1]{\textbf{\textcolor{blue}{#1}}}
\newcommand{\lex}[1]{\textbf{\mtcitestyle{#1}}}


% \renewcommand{\labelitemi}{\textcolor{violet}{‣}}
% \renewcommand{\labelitemii}{\textcolor{purple}{‣}}

\newcommand{\subhead}[1]{\noindent\textbf{#1}\\[5mm]}

\newcommand{\Bad}{\emp{\raisebox{0.15ex}{\ensuremath{\mathbf{\otimes}}}}}
\newcommand{\bad}[1]{*\eng{#1}}

\newcommand{\com}[1]{\hfill (\emp{#1})}%

\usepackage{relsize,xspace}
\newcommand{\into}{\ensuremath{\rightarrow}\xspace}
\newcommand{\tot}{\ensuremath{\leftrightarrow}\xspace}
\usepackage{url}
\newcommand{\lurl}[1]{\MyLogo{\url{#1}}}

\usepackage{langsci-gb4e}
\let\eachwordone=\slshape
%%% langsci-gb4e defines \thexnumi as \@roman\@xsi{xnumi}, which only works
%%% inside an exe environment; some decks use xlisti on its own.  Restore
%%% the plain gb4e definition so that keeps working.
\makeatletter
\def\thexnumi{\@xsi{xnumi}}
\makeatother
\newcommand{\lx}[1]{\textbf{\mtciteform{#1}}}
\newcommand{\ix}{\ex\slshape}



\newcommand{\ent}{\ensuremath{\Rightarrow}\xspace}
\newcommand{\ngv}{\ensuremath{\not\Rightarrow}\xspace}
%\usepackage{times}
%\usepackage{nttfoilhead}
\newcommand{\myslide}[1]{\foilhead[-25mm]{\raisebox{12mm}[0mm]{\red{#1}}}\MyLogo{\logo}}
\newcommand{\myslider}[1]{\rotatefoilhead[-25mm]{\raisebox{12mm}[0mm]{\red{#1}}}}
%\newcommand{\myslider}[1]{\rotatefoilhead{\raisebox{-8mm}{\emp{#1}}}}

\newcommand{\section}[1]{\myslide{}{\begin{center}\Huge \red{#1}\end{center}}}



\usepackage[lyons,j,e,k]{mtg2e}
%\renewcommand{\mtcitestyle}[1]{\textcolor{-red!75!green!50}{\textsl{#1}}}
\renewcommand{\mtcitestyle}[1]{\textcolor{teal}{\textsl{#1}}}
\newcommand{\cs}{\mtciteform}
\newcommand{\iz}[1]{\texttt{\textup{#1}}}
\newcommand{\gm}{\textsc}
\usepackage[normalem]{ulem}
\newcommand{\ul}{\uline}
\newcommand{\ull}{\uuline}
\newcommand{\wl}{\uwave}
\newcommand{\vs}{\ensuremath{\Leftrightarrow}~}
%%%
%%% Bibliography
%%%
\usepackage{natbib}
%\usepackage{url}
\usepackage{bibentry}



\usepackage[utf8]{inputenc}
%%% trees
\usepackage{multicol}
%%% pst-node: one example in lec-12-control draws a connector between
%%% two points with \rnode/\ncbar.  This used to come in with rtrees.
\usepackage{pst-node}
\usepackage{forest}
\useforestlibrary{linguistics}
\forestapplylibrarydefaults{linguistics}

%%% \makebox's optional arguments are read as AVM delimiters by
%%% langsci-avm, so hide them behind a macro with no visible brackets.
\newcommand{\lbox}[1]{\makebox[0mm][l]{#1}}

%%% Tree styles.
%%%   words  -- trees whose leaves are actual words (italicised, aligned)
%%%   cat    -- trees whose leaves are categories or feature structures
\forestset{
  words/.style={for tree={s sep=7mm, l sep=5mm,
                          if n children=0{font=\itshape, tier=terminal}{}}},
  cat/.style={for tree={s sep=7mm, l sep=5mm}},
}
\newcommand{\wtree}[1]{\begin{forest} words #1 \end{forest}}
\newcommand{\ctree}[1]{\begin{forest} cat #1 \end{forest}}
%%% the rule being applied at this step of a derivation
\newcommand{\step}[1]{\begin{center}\emp{#1}\end{center}}

%%% AVMs (Language Science Press)
\usepackage{langsci-avm}

%%% Macros the decks use, carried over from headx.tex
\newcommand{\blank}{\rule{3em}{1pt}\xspace}
\newcommand{\ft}[1]{\textsc{#1}}
\newcommand{\idn}{\mtciteform}    %%% Indonesian
\newcommand{\prd}[1]{\textbf{#1}}
\newcommand{\typ}[1]{\textit{#1}}
%%% \val: provide before renewing, in case a package defines it.
\providecommand{\val}[1]{}
\renewcommand{\val}[1]{\textit{#1}}
%%% Manning's \avmbox, still used in tree labels and in maths where a
%%% bare langsci \tag would not work.  \avm{\tag{...}} gives the same
%%% boxed tag as the ones inside feature structures.
\newcommand{\avmbox}[1]{\mbox{\avm{\tag{#1}}}}

%%% Ring a value in a dotted oval.  These used pstricks \ovalnode,
%%% which came in with rtrees; tikz (via forest) does the same job.
\newcommand{\OV}[1]{\tikz[baseline]{\node[draw,dotted,red,ellipse,
          inner sep=1pt,outer sep=0pt]{#1};}}
\newcommand{\OVB}[1]{\tikz[baseline]{\node[draw,dotted,blue,ellipse,
          inner sep=1pt,outer sep=0pt]{#1};}}

%%% From CSLI book
\newcommand{\mc}{\multicolumn}
\newcommand{\HD}{\textbf{H}\xspace}
\newcommand{\el}{< >}       %%% empty list; langsci-avm needs the space

\newcommand{\A}{\noindent\textbf{A}: }
\newcommand{\Q}{\noindent\textbf{Q}: }

%%% Question and answer on two pages, written once:
%%%   \qaslide{Title}{... question ... \answer{... answer ...}}
%%% prints the slide twice, first without the \answer parts and then
%%% with them, so the answer appears "on click".  The example counter
%%% is reset so both pages show the same example numbers: gb4e keeps the
%%% running number in the equation counter between exe environments, and
%%% in xnumi inside one.
\newif\ifshowanswer
\newcommand{\answer}[1]{\ifshowanswer#1\fi}
\newcounter{qaxnumi}
\newcounter{qaeqn}
\newcommand{\qaslide}[2]{%
  \setcounter{qaxnumi}{\value{xnumi}}%
  \setcounter{qaeqn}{\value{equation}}%
  \showanswerfalse\myslide{#1}#2%
  \setcounter{xnumi}{\value{qaxnumi}}%
  \setcounter{equation}{\value{qaeqn}}%
  \showanswertrue\myslide{#1}#2%
  \showanswerfalse}
%\newcommand{\C}{\noindent\textbf{C}: }


\usepackage{tikz}
\usepackage{tikz-qtree}
\newcommand{\bx}[1]{\avmbox{\smaller #1}}

\begin{document}

\header{Lecture 14}{Long Distance Dependencies}{}
\maketitle


\overview
\MyLogo{\citet{sag:wasow:bender:2003} --- Chapter 14}

% \section{Review of Binding}

% \myslide{Binding deals with pronouns}

% \begin{exe}
% \ix She likes herself
% \ix *She$_i$ likes her$_i$.
% \ix We gave presents to ourselves.
% \ix *We gave presents to us.
% \ix We gave ourselves presents
% \ix *We gave us presents.
% \ix *Leslie told us about us.
% \ix Leslie told us about ourselves.
% \ix *Leslie told ourselves about us.
% \ix *Leslie told ourselves about ourselves.
% \end{exe}

% \myslide{Some Terminology}
% \begin{itemize}
% \item \txx{Binding}: The association between a pronoun
% and an antecedent.
% \item \txx{Anaphoric}: A term to describe an element (e.g.
% a pronoun) that derives its interpretation from
% some other expression in the discourse.
% \item \txx{Antecedent}: The expression an anaphoric
% expression derives its interpretation from.
% \item \txx{Anaphora}: The relationship between an
% anaphoric expression and its antecedent.
% \end{itemize}

% \myslide{The Argument Realization Principle}
% \begin{itemize}
% \item We introduce a feature ARG-ST:\\
% \begin{avm}\avmfont{\sc} 
% \[ \asort{word}  
%  syn & \[ val & \[ spr &  \@{A} \\comps &  \@{B} \] \] \\
%  arg-st & \< \@{A} $\oplus$ \@{B} \> \]
% \end{avm}
% \item This is a constraint on the type \val{word}
% \begin{itemize}
% \item It’s neither in \ft{syn} nor \ft{sem}.
% \item It only appears on lexical heads (not
% appropriate for type phrase)
% \item No principle stipulates identity
% between \ft{arg-st}s.
% \end{itemize}
% \end{itemize}

% \myslide{The Binding Principles}
% \begin{itemize}
% \item \txx{Principle A}: A [\ft{mode} \val{ana}] element must be
% outranked by a coindexed element.
% \item \txx{Principle B}: A [\ft{mode} \val{ref}] element must not
% be outranked by a coindexed element.
% \item \txx{Anaphoric Agreement Principle} (AAP):
% Coindexed NPs agree.
% \end{itemize}

% \begin{description}
% \item [Formalization] ~\\[-2ex]
% \begin{itemize}
% \item Definition: If A precedes B on some \ft{arg-st} list,
% then A \txx{outranks} B.
% \item Elements that must be anaphoric --- that is, that
% require an antecedent --- are lexically marked
% [\ft{mode} \val{ana}]. These include reflexive pronouns
% and reciprocals.
% \end{itemize}

% \end{description}
% \myslide{Binding in PPs}
% \begin{itemize}
% \item The Binding Principles by themselves don’t block:
%   \begin{exe}
%     \ix I brought a book with me.
%     \ix *I brought a book with myself.
%     \ix *I mailed a book to me.
%     \ix I mailed a book to myself.
%   \end{exe}
% \item Two Types of Prepositions
% \begin{itemize}
% \item \txx{Argument-marking}: Function like casemarkers, indicating the
% roles of NP referents in the situation denoted by the verb.
%  \begin{itemize}
%   \item Share their objects' \ft{mode} and \ft{index} values.
%   \end{itemize}
% \item \txx{Predicative}: Introduce their own predication.
%   \begin{itemize}
%   \item Introduce their own \ft{mode} and \ft{index} values.
% \end{itemize}
% \end{itemize}
% \end{itemize}


% % \myslide{Redefining Rank}
% %  \begin{itemize}
% % \item If there is an \ft{arg-st} list on which A
% % precedes B, then A outranks B.
% % \item If a node is coindexed with its daughter, they
% % are of equal rank -- that is, they outrank the
% % same nodes and are outranked by the same
% % nodes.
% % \end{itemize}

% \myslide{\eng{I sent a letter to myself}}

% \begin{avmtree}\avmfont{\sc} 
% \br{S}{
%   \br{\@1 NP$_i$}{ \lf{\eng{I}} }
%   \br{VP \[ \q< \@1 \q> \]}{ 
%     \br{V \[ spr & \q< \@1 \q> \\
%              comps & \q< \@2, \@3 \q> \\ 
%              arg-st & \q< \@1, \@2, \@3 \q> \]}{
%       \lf{\eng{sent}}}
%     \br{\@2 NP$_j$}{
%       \br{D}{\lf{\eng{a}}}
%       \br{N$_j$}{\lf{\eng{letter}}}}
%     \br{\@3 PP$_i$}{
%       \br{P$_i$}{\lf{\eng{to}}}
%       \br{NP$_i$ \[ mode & ana \]}{\lf{\eng{myself}}}}
% }}
% \end{avmtree}



% \myslide{The Imperative Rule}

% \begin{avm}\avmfont{\sc} 
% \[ \asort{phrase}  
%    head & verb \\
%    val & \[ spr  & \<  \> \] \\
%    sem & \[ mode & dir \\ index & $s$ \] \]
%               \ $\rightarrow$\ \HD \ 
% \[ \asort{phrase}  
%    head & \[ \asort{verb} form & base \]\\
%    val & \[ spr  & \< NP \[ per 2nd \]  \> \] \\
%    sem & \[ index & $s$ \] \]
% \end{avm}


% \begin{itemize}\addtolength{\itemsep}{-2ex}
% \item Internal structure of a VP
% \item Directive function
% \item Base form
% \item Only 2nd person reflexives
% \end{itemize}

% \myslide{Imperative example}
% %\begin{multicols}{2}
% \begin{avmtree}\avmfont{\sc} 
% \br{S}{
%   \br{VP \[ spr & \q< \@1 NP$_h$ \[ per 2nd \\ num sg \] \]}{ 
%     \br{V \[ spr & \q< \@1 \q>  \\ arg-st & \q< \@1 NP$_h$,  \@2 PP$_i$ \q> \]}{
%       \lf{\eng{vote}}}
%     \br{\@2 PP$_i$}{
%       \br{P$_i$}{\lf{\eng{for}}}
%       \br{NP$_i$ \[ mode & ana \]}{\lf{\eng{yourself}}}}
% }}
% \end{avmtree}
% %\newpage

% % \begin{itemize}
% % \item What's the SPR value on S?
% % %Why?
% % \item What's the SPR value on VP?
% % %Why?
% % \item What's the SPR value on V?
% % %Why?
% % \item Which nodes have ARG-ST?
% % \item Which ARG-ST matters for
% % the licensing of \eng{yourself}?
% % \end{itemize}
% % %\newpage

\qaslide{A Note on Adjectives}{

\begin{itemize}\addtolength{\itemsep}{-1.5ex}
 \item Attributive adjectives are related to predicative adjectives by
   a lexical rule that coindexes the first element of \ft{arg-st} with
   \ft{mod} and sets \ft{spr} to an empty list: 
\\   \eng{The dog is red}  \into \eng{The red dog}
   \\    \scalebox{0.8}{\avm{< \textnormal{red}, \ [\type*{word}
      head &  [\type*{adj} pred & $+$ ]  \\
      % [\type*{agr-cat} agr & [ per & 3 \\ num & sg ] ]
      val & [ spr & < \1 NP >\\
              comps & <  >\\
              mod & <  > ] \\
      arg-st &      < \1  >   
      ]> \ \
\into \ \ 
  < \textnormal{red}, \ [\type*{word}
      head & \type{adj} \\
      % [\type*{agr-cat} agr & [ per & 3 \\ num & sg ] ]
      val & [ spr & <  >\\
              comps & <  >\\
              mod & < \1 NP > ] \\
  arg-st &      < \1  >   
     ]>
}}
 \item We can't just coindex \ft{spr} and \ft{mod}.  Why?
  \answer{\\ {\small\emp{Then it would still want a subject, but attributive adjectives take none.}}}
 \item The \ft{spr} adjectives need a
   subject-raising \eng{be} to form a sentence. Why?
  \answer{\\ {\small\emp{Only (finite) verbs head sentences: \lex{be} is the verb, and raises the subject.}}}
 \end{itemize}
}

\section[Long-distance dependencies: some examples, and what is new and different about them]{Long-Distance Dependencies}


\myslide{Examples}
\begin{multicols}{2}
\begin{itemize}
\item \emp{Grammatical}:
  \begin{exe}
    \ex \eng{Did you find something?}  
    \ex \eng{Tell me you talked to someone!}
  \end{exe}
\item \txx{wh-questions}:
  \begin{exe}
    \ex \eng{What did you find? } 
    \ex \eng{Tell me who you talked to.}
  \end{exe}
\item  \txx{relative clauses}:
  \begin{exe}
    \ex \eng{the item that I found}
    \ex \eng{the guy who(m) I talked to}
  \end{exe}
\item \emp{Ungrammatical}:
  \begin{exe}
    \ex \bad{did you find}
    \ex \bad{you talked to}
  \end{exe}
\item  \txx{topicalization}:
  \begin{exe}
    \ex \eng{The manual, I can't find.}
    \ex \eng{Chris, you should talk to.}
  \end{exe}
\item  \lex{easy}-\txx{adjectives}:
  \begin{exe}
    \ex \eng{My house is easy to find.}
    \ex \eng{Pat is hard to talk to.}
  \end{exe}
\end{itemize}
\end{multicols}

\myslide{What These Have in Common}
\begin{itemize}
\item There is a \txx{gap}: nothing following \lex{find} and \lex{to},
even though both normally require objects.
\item Something that fills the role of the element
missing from the gap occurs at the beginning of
the clause.
\item We use topicalization and \lex{easy}-adjectives to
illustrate the phenomenon:
\begin{exe}
  \ex \eng{The manual$_g$, I can't find \blank$_g$.}
  \ex \eng{Chris$_g$ is easy to talk to \blank$_g$.}
\end{exe}
\end{itemize}

\myslide{Gaps and Their Fillers Can Be Far Apart}
\begin{exe}
  \ex \eng{\ul{The solution to this problem}$_g$, Pat said that
  someone claimed you thought I would never
  find \blank$_g$.}
  \ex \eng{Chris$_g$ is easy to consider it impossible for anyone
  but a genius to try to talk to \blank$_g$.}
\end{exe}
\vspace*{-2ex}
\begin{itemize} 
\item Fillers often have syntactic properties associated with their
  gaps
\end{itemize}
\begin{exe}
\ex
\begin{xlist}
  \ex \eng{Him$_g$, I haven’t met \blank$_g$.}
  \ex \bad{He$_g$, I haven’t met \blank$_g$.}
\end{xlist}
\ex
\begin{xlist}
  \ex \eng{The scissors$_g$, Pat told us \blank$_g$ were missing.}
  \ex \bad{The scissors$_g$, Pat told us \blank$_g$ was missing.}
\end{xlist}
\ex
\begin{xlist}
  \ex \eng{\ul{On Pat}$_g$, you can rely \blank$_g$.}
  \ex \bad{\ul{To Pat}$_g$, you can rely \blank$_g$.}
\end{xlist}
\end{exe}
\vspace*{-3ex}
\begin{itemize}
\item That’s why we call them \txx{long-distance \ul{dependencies}}
\end{itemize}


% \myslide{LDDs in TG}
% \begin{itemize}
% \item These were long thought to constitute the
% strongest evidence for transformations.
% \item They were handled in TG by moving the filler
% from the gap position.
% \item Case, agreement, preposition selection could
% apply before movement.
% \end{itemize}

% \myslide{A big debate about LDDs in TG}
% \begin{itemize}
% \item Does long-distance movement take place in one fell swoop
% or in lots of little steps?
% \end{itemize}
% % ↑

% % ↑

% % ↑
% % ↑

% Swooping

% Looping


\myslide{Other Relevant Facts}

\begin{itemize}
\item Various languages show morphological marking
on the verbs or complementizers of clauses
between the filler and the gap.
\item Psycholinguistic evidence indicates increased
processing load in the region between filler and
gap.
% \item This opens the door to non-transformational
% analyses, in which the filler-gap dependency is
% mediated by local information passing.
\item So something is being kept track of all the way from the gap
  to the filler.  That is what our analysis will do: pass the
  information up step by step, through local information passing.
\end{itemize}

\section{Our Approach: The Feature \ft{gap}}

\myslide{A Rough Sketch of Our Approach}
\begin{itemize}
\item A feature \ft{gap} records information about a
missing constituent.
\item The \ft{gap} value is passed up the tree by a new
principle.
\item A new grammar rule allows us to expand S as a filler
followed by another S whose \ft{gap} value
matches the filler.
\item Caveat: Making the details of this general
idea work involves several complications.
\item The core idea comes from \citet{Gazdar:1981}.
\end{itemize}



\myslide{The Feature \ft{gap}}
\begin{itemize}
\item Like valence features and \ft{arg-st}, \ft{gap}’s
value is a list of feature structures (often
empty).  You can have multiple gaps.
\item Subject gaps are introduced by a lexical rule.
\item Non-subject gaps are introduced by revising
the Argument Realization Principle.
\end{itemize}

\myslide{The Revised ARP}

\avm{[\type*{word}  
 syn & [ val & [ spr &  \tag{A} \\comps &  \tag{B} $\ominus$ \tag{C} ] \\
          gap & \tag{C} ] \\
 arg-st & \tag{A} $\oplus$ \tag{B} ]
}
\begin{itemize}
\item $\ominus$ is a kind of list subtraction 
  \begin{itemize}
  \item It's not always defined (the sublist must exist on the main list)
  \item When defined, it's not always unique 
  \end{itemize}
\item The ARP now says the non-\ft{spr} arguments are
distributed between \ft{comps} and \ft{gap}.
\end{itemize}

\qaslide{List Subtraction}{
\begin{itemize}\addtolength{\itemsep}{-0.5ex}
\item $\langle$ NP, PP $\rangle \ominus \langle$ NP $\rangle$ = ?
  \answer{\emp{$\langle$ PP $\rangle$}}
\item $\langle$ NP, PP, NP $\rangle \ominus \langle$ NP $\rangle$ = ?
  \answer{\emp{$\langle$ PP, NP $\rangle$ or $\langle$ NP, PP $\rangle$: not unique}}
\item $\langle$ NP, PP, NP $\rangle \ominus \langle$ NP, PP $\rangle$ = ?
  \answer{\emp{$\langle$ NP $\rangle$}}
\item $\langle$ NP, PP, NP $\rangle \ominus \langle$ NP, VP $\rangle$ = ?
  \answer{\emp{undefined: there is no VP to take away}}
\item $\langle$ VP, PP, NP $\rangle \ominus \langle$ NP, VP $\rangle$ = ?
  \answer{\emp{undefined: order matters, the VP is not after the NP}}
\end{itemize}
\answer{\emp{So for \lex{hand} with \ft{arg-st} $\langle$ NP, NP, PP
    $\rangle$ the ARP allows three words: nothing in \ft{gap}, the
    NP in \ft{gap}, or the PP in \ft{gap} (next slides).}}
}

\myslide{A Word with a Non-Empty \ft{gap} Value}

    \avm{< \textnormal{hand}, \ [\type*{word}
      syn & [ head & [\type*{verb} form & fin ] \\
      % [\type*{agr-cat} agr & [ per & 3 \\ num & sg ] ]
      val & [ spr & < \1 NP [ case & nom \\ agr & non-3sing ] > \\ 
      comps & < \3 PP:to > ] \\
      gap &  < \tag{2NP} [ case & acc ] > ] \\
      arg-st & < \1, \2, \3 >
      ]>
}

    \myslide{A Word with Another Non-Empty \ft{gap} Value}

    \avm{< \textnormal{hand}, \ [\type*{word}
      syn & [ head & [\type*{verb} form & fin ] \\
      % [\type*{agr-cat} agr & [ per & 3 \\ num & sg ] ]
      val & [ spr & < \1 NP [ case & nom \\ agr & non-3sing ] > \\ 
      comps & < \tag{2NP} [ case & acc ] > ] \\
      gap &  <  \3 PP:to > ] \\
      arg-st & < \1, \2, \3 >
      ]>
}

\myslide{The Same Word with an Empty \ft{gap} Value}
\MyLogo{The ARP gives rise to multiple lexemes}

    \avm{< \textnormal{hand}, \ [\type*{word}
      syn & [ head & [\type*{verb} form & fin ] \\
      % [\type*{agr-cat} agr & [ per & 3 \\ num & sg ] ]
      val & [ spr & < \1 NP [ case & nom \\ agr & non-3sing ] > \\ 
      comps & < \tag{2NP} [ case & acc ], \3 PP:to > ] \\
      gap &  <  > ] \\
      arg-st & < \1, \2, \3 >
      ]>
}

    
% \myslide{How We Want \ft{gap} to Propagate}
% \MyLogo{Traditionally called \ft{slash}: VP/NP}
% \begin{avmtree}\avmfont{\sc}
% \br{S \[ gap & < > \]}{ \br{NP \[ gap & < > \]}{\lf{Kim}}
%   \br{S/NP \[ gap & \q< NP \q> \]}{\br{NP \[ gap & < > \]}{\lf{we}}
%     \br{VP/NP \[ gap & \q< NP \q> \]}{\br{V \[ gap & < > \]}{\lf{know}}
%       \br{S/NP \[ gap & \q< NP \q> \]}{\br{NP \[ gap & < > \]}{\lf{Bobby}}
%         \br{VP/NP \[ gap & \q< NP \q> \]}{\lf{hates}}}}}}
% \end{avmtree}


\myslide{How We Want \ft{gap} to Propagate}
\MyLogo{Traditionally called \ft{slash}: VP/NP
  (think of it as a VP missing an NP)}
% Kim we know bobby likes
\newcommand{\avn}[2]{#1\\\smaller\avm{[ #2 ]
}} %avm node
\newcommand{\avN}[2]{#1\,\smaller\avm{[ #2 ]
}} %avm node
\newcommand{\avX}[1]{\avm{[ #1 ]
}} %avm node

\scalebox{0.8}{\begin{tikzpicture}
\tikzset{level distance=8ex}
\tikzset{frontier/.style={distance from root=40ex}}
\tikzset{every tree node/.style={align=center,anchor=north}}
\Tree [.\avn{S}{gap & < >} 
         [.\avn{NP}{gap & < >} Kim ]
         \edge[style=double,very thick];
         [.\avn{S/NP}{gap & < NP >} 
           [.\avn{NP}{gap & < >} we ]
           \edge[style=double,very thick];
           [.\avn{VP/NP}{gap & < NP >} 
           \edge[style=double,very thick];
              [.\avn{V}{gap & < >} know ]
              [.\avn{S/NP}{gap & < NP >} 
                [.\avn{NP}{gap & < >} Bobby ]
                \edge[style=double,very thick];
                [.\avn{VP/NP}{gap & < NP >} likes ] ] ] ] ]
\end{tikzpicture}}
                              

\section[Filling gaps: \ft{stop-gap}, the Head-Filler Rule and \lex{easy}-adjectives]{Filling Gaps}

\myslide{What \ft{gap} Propagation Should Be Doing}
\MyLogo{}
\begin{itemize}
\item Pass any \ft{gap} values from daughters up to their
mothers,
\item[\ldots] \emp{except} when the filler is found.
\item For topicalization, we can write the exception into
the grammar rule.
\item For \lex{easy}-adjectives, the NP that corresponds to the
gap is the subject, which is introduced by the
Head-Specifier Rule.
\item Since specifiers are not generally gap fillers, we
can't write the gap-filling into the HSR.
\end{itemize}

\myslide{Our Solution to This Problem}
\begin{itemize}
\item For \lex{easy}-adjectives, we treat the adjective formally
as the filler, marking its \ft{spr} value as coindexed
with its \ft{gap} value.
\item We use a feature \ft{stop-gap} to trigger the
  emptying of the \ft{gap} list.
  \begin{itemize}
  \item \ft{stop-gap} stops gap propagation
  \item \lex{easy}-adjectives mark \ft{stop-gap} lexically
  \item A new grammar rule, the \txx{Head-Filler Rule},
    contains \ft{stop-gap}
  \end{itemize}
\end{itemize}

\myslide{The \ft{gap} Principle}
A local subtree $\Phi$ satisfies the \ft{gap} Principle with respect to a
headed rule  if and only if $\Phi$ satisfies:


% \begin{avmtree}\avmfont{\sc}
%      \br{\[ gap & \( \@{A_1} \ $\oplus$\  \ldots \ $\oplus$\ \@{A_n} \)
%        \ $\ominus$\  \@{A_0}\]}{
%      \br{\[ gap & \@{A_1} \]}{\tlf{~}}
%       \lf{\ldots}
%       \br{\HD \[ gap & \@{A_i} \\ stop-gap & \@{A_0} \]}{\tlf{~}}
%       \lf{\ldots}
%       \br{\[ gap & \@{A_n} \]}{\tlf{~}}}
% \end{avmtree}

\begin{tikzpicture}
\tikzset{level distance=8ex}
   \Tree [.\avX{gap & ( \tag{$A_1$} \ $\oplus$\  \ldots \ $\oplus$\
                         \tag{$A_n$} ) $\ominus$\  \tag{$A_0$}} 
            [.\avX{gap & \tag{$A_1$}} ]
            [.{\ldots}  ]
            \edge[style=double,very thick];
            [.\avN{H}{gap & \tag{$A_i$}\\ stop-gap & \tag{$A_0$}}   ]
            [.{\ldots}  ]
            [.\avX{gap & \tag{$A_n$}} ] ]
\end{tikzpicture}

\begin{itemize}
\item The \ft{gap} of the mother is the append of the \ft{gap}s of the daughters
\item [\ldots{}] minus \ft{stop-gap} on the head daughter
 
\end{itemize}


\myslide{How Does \ft{stop-gap} Work?}
\begin{itemize}
\item \ft{stop-gap} is empty almost everywhere
\item When a gap is filled, \ft{stop-gap} is nonempty,
and its value is the same as the gap being filled.
\item This blocks propagation of that \ft{gap} value, so
gaps are only filled once.
\item The nonempty \ft{stop-gap} values come from two
  sources:
  \begin{itemize}
  \item A stipulation in the Head-Filler Rule
  \item Lexical entries for \lex{easy}-adjectives
  \end{itemize}
\item No principle propagates \ft{stop-gap}
\end{itemize}

\myslide{The Head-Filler Rule}

\avm{[\type*{phrase}] \ \ \into \ \ 
\ \ \1 [ gap & < > ]\ \  
\HD [ head & [\type*{verb} form & fin ] \\
       val & [ spr & < >\\
                comps & < >] \\
       gap & < \1 > \\
       stop-gap & < \1 > ]
}
\begin{itemize}\addtolength{\itemsep}{-2ex}
\item This only covers gap filling in finite Ss
\item The filler has to be identical to the \ft{gap} value
\item The \ft{stop-gap} value is also identical
\item The \ft{gap} Principle ensures that the mother's \ft{gap} value is the
empty list
\end{itemize}

\myslide{Gap Filling with \lex{easy}-Adjectives}
  \begin{exe}
    \ex \label{l:easy}
 \scalebox{0.8}{\avm{< \textnormal{easy}, \ [\type*{word}
      syn & [ head & \type{adj} \\
      % [\type*{agr-cat} agr & [ per & 3 \\ num & sg ] ]
      val & [ spr & < \1 >\\
              comps & < \3 > ] \\
      stop-gap & < \2 >] \\
      arg-st &      < \1 NP$_i$,  \3 VP [ inf & + \\ gap & < \2 NP$_i$, \ldots > ] >   
      ]> \ \
}}

  \end{exe}
\begin{itemize}\addtolength{\itemsep}{-1ex}
\item Because \ft{stop-gap} and \ft{gap} have the same value, that
value will be subtracted from the mother's \ft{gap} value.
\item The first argument is coindexed with the \ft{gap} value,
accounting for the interpretation of the subject as the filler.
\end{itemize}

\myslide{Other \lex{easy}-Adjectives}
\begin{center}\large
  \lex{alarming}, \lex{boring}, \lex{convenient}, \lex{cheap},
  \lex{dangerous}, \lex{entertaining}, \lex{fruitless}, \lex{gross},
  \lex{hazardous}, \lex{illegal}, \lex{kind}, \lex{logical},
  \lex{motivating}, \lex{natural}, \lex{optional}, \lex{prudent},
  \lex{risky}, \lex{sobering}, \lex{tedious}, \lex{useful},
  \lex{valuable}, \lex{wise}, \ldots
\end{center}
\begin{itemize}
\item All take an infinitival VP with a gap whose filler is the
  subject: \eng{This form is tedious to fill in \blank},
  \eng{Pat is dangerous to argue with \blank}.
\item Same entry as \lex{easy} with a different predication.
\end{itemize}

\qaslide{Why Do We Need \ft{stop-gap} at All?}{
\begin{itemize}
\item Why not treat \ft{gap} like the valence features: give the
  Head-Filler Rule a [\ft{gap} \el] mother and let the non-head
  daughter be the filler?
  \answer{\\ {\small\emp{That would work for topicalization, where the
    filler is a daughter of the rule.  But \lex{easy}-adjectives do not
    occur in the Head-Filler Rule: \eng{easy to talk to} is an ordinary
    Head-Complement phrase, and \eng{Kim is easy to talk to} ordinary
    Head-Specifier phrases.  The \ft{gap} Principle would pass the
    complement's gap up through them, and no rule would be there to
    cancel it.  \ft{stop-gap} lets the \emph{word} \lex{easy} do the
    cancelling, wherever it occurs.}}}
\end{itemize}
}

\myslide{A Tree for \eng{easy to talk to \blank}}
\begin{tikzpicture}
\tikzset{level distance=13ex}
%\tikzset{frontier/.style={distance from root=40ex}}
\tikzset{level 2/.style={level distance=20ex}}
\tikzset{every tree node/.style={align=center,anchor=north}}
\Tree [.\avn{AP}{val & [ spr & < \tag{2NP}$_i$ > ]\\ gap & < >}
\edge[style=double,very thick];
         [.\avn{A}{ val & [ spr & < \2 > \\ comps & < \3 > ]\\ 
                    gap & < ~ > \\ 
                    stop-gap & < \1 >} easy ]
         [.\avn{\bx{3} VP/NP}{val & [ spr & < NP > ]\\ 
                           gap & < \1 NP$_i$ >} \edge[roof]; {to talk to} ] ]

\end{tikzpicture}
% \begin{avmtree}\avmfont{\sc}

% \br{AP \[ val & \[ spr & \< \@2NP$_i$ \> \]\\ gap & < > \]}{
%   \br{A \[ val & \[ spr & \< \@2 \> \\ comps & \< \@3 \> \]\\ 
%     gap & \< ~ \> \\ stop-gap & \< \@1 \>\]}{\lf{easy}}
%   \br{\@3 VP\[ val & \[ spr & \< NP \> \]\\ 
%                gap & \< \@1 NP$_i$ \> \]}{\tlf{to talk to \blank}}}

% \end{avmtree}

\myslide{\ft{stop-gap} Housekeeping}
\begin{itemize}
\item Lexical entries with nonempty \ft{stop-gap}
values (like \lex{easy}) are rare, so \ft{stop-gap} is by
default empty in the lexicon.
\item Head-Specifier and Head-Modifier Rules need to
say [\ft{stop-gap} $<$ $>$].
\item Lexical rules preserve \ft{stop-gap} values.
\end{itemize}

\myslide{Where Is \ft{stop-gap} Specified?}
\begin{center}\small
\begin{tabular}{p{0.36\textwidth}p{0.52\textwidth}}
\textbf{Where} & \textbf{Value} \\ \hline
lexemes & default \avm{[ stop-gap & / < > ]}; \lex{easy}-adjectives override it \\[0.5ex]
head daughter of the Head-Complement Rule & a word, so the lexical default decides \\[0.5ex]
head daughter of the Head-Specifier and Head-Modifier Rules & the rules say \avm{[ stop-gap & < > ]} \\[0.5ex]
head daughter of the Head-Filler Rule & the rule says \avm{[ stop-gap & < \1 > ]}, the filler \\[0.5ex]
mothers of any rule & unspecified; fixed when the phrase becomes a head daughter higher up \\
\end{tabular}
\end{center}
\begin{itemize}
\item No principle passes \ft{stop-gap} anywhere: it is set where it is
  used and nowhere else.
\end{itemize}

\qaslide{\ft{gap} Housekeeping}{
\begin{itemize}
\item The initial symbol must say [\ft{gap} $< >$ ]. Why?
  \answer{\\ {\small\emp{To block \bad{Pat found} and \bad{Chris talked to} as
    stand-alone sentences.}}}
\item The Imperative Rule must propagate \ft{gap} values.  Why?
  \answer{\\ {\small\emp{It's not a headed rule, so the effect of the \ft{gap}
    Principle must be replicated, and imperatives can have gaps:
    \eng{This book, put on the top shelf!}}}}
\end{itemize}
}

\myslide{Sentences with Multiple Gaps}
\begin{itemize}
\item Famous examples:
  \begin{exe}
    \ex \eng{This violin$_i$, sonatas$_j$ are easy to play \blank$_j$ on \blank$_i$.}
    \ex ?\eng{Sonatas$_j$, this violin$_i$ is easy to play \blank$_j$ on \blank$_i$.}
  \end{exe}
  The second is licensed, but only with the odd reading in which the
  violin is played on the sonatas.
\item Our analysis gets this:
  \begin{itemize}
  \item The subject of \lex{easy} is coindexed with the first
    element of the \ft{gap} list.
  \item The Head-Filler Rule only allows one \ft{gap}
    remaining.
  \item Step by step on the next slides.
  \end{itemize}
\item There are languages that allow multiple gaps more
generally.
\end{itemize}

\qaslide{Step by Step: \ft{gap} and \ft{stop-gap} on Every Node}{
\MyLogo{\eng{This violin, sonatas are easy to play \blank\ on \blank}}
\begin{minipage}[t]{0.6\textwidth}\vspace{0pt}
\scalebox{0.52}{\begin{tikzpicture}
\tikzset{level distance=8.5ex}
\tikzset{every tree node/.style={align=center,anchor=north}}
\Tree [.\avn{S}{gap & < >}
        [.\avn{\bx{1}NP$_i$}{gap & < >} \edge[roof]; {This violin} ]
        \edge[style=double,very thick];
        [.\avn{S}{gap & < \1 > \\ stop-gap & < \1 >}
          [.\avn{\bx{2}NP$_j$}{gap & < >} sonatas ]
          \edge[style=double,very thick];
          [.\avn{VP}{gap & < \1 >}
            \edge[style=double,very thick];
            [.\avn{V}{gap & < >} are ]
            [.\avn{AP}{gap & < \1 >}
              \edge[style=double,very thick];
              [.\avn{A}{gap & < > \\ stop-gap & < \2 >} easy ]
              [.\avn{VP}{gap & < \2, \1 >}
                \edge[style=double,very thick];
                [.\avn{V}{gap & < >} to ]
                [.\avn{VP}{gap & < \2, \1 >}
                  \edge[style=double,very thick];
                  [.\avn{V}{gap & < \2 NP$_j$ >} play ]
                  [.\avn{PP}{gap & < \1 NP$_i$ >}
                    \edge[style=double,very thick];
                    [.\avn{P}{gap & < \1 NP$_i$ >} on ] ] ] ] ] ] ] ]
\end{tikzpicture}}
\end{minipage}\hfill
\begin{minipage}[t]{0.38\textwidth}\vspace{0pt}
\begin{itemize}
\item Fill in \ft{gap} on every node, and \ft{stop-gap} where it is
  not empty.
\end{itemize}
\answer{{\small\emp{\eng{play} and \eng{on} each put their object in
  \ft{gap}; the Head-Complement Rule appends them, V first: $\langle$
  \bx{2}, \bx{1} $\rangle$.  \lex{easy} cancels the first and
  coindexes it with its subject \eng{sonatas}; the Head-Filler Rule
  fills the one that is left with \eng{this violin}.  Swap the
  fillers and the violin is what gets played.}}}
\end{minipage}
}

\myslide{Where We Are}

\begin{itemize}
\item \txx{filler-gap} structures:
  \begin{exe}
   \ex \eng{\ul{The solution to this problem}, nobody
   understood \blank.}
   \ex \eng{\ul{That problem} is easy to understand \blank.}
 \end{exe}
\item The feature \ft{gap} encodes information about
missing constituents
\item Modified ARP allows arguments that should be on
the \ft{comps} list to show up in the \ft{gap} list
\item \ft{gap} values are passed up the tree by the \ft{gap}
Principle
\end{itemize}

\myslide{Where We Are (Continued)}
\begin{itemize}
\item The feature \ft{stop-gap} signals where \ft{gap} passing
should stop
\item The Head-Filler Rule matches a filler to a \ft{gap} and
(via \ft{stop-gap}) empties \ft{gap}
\item Lexical entries for \lex{easy}-adjectives require a gap in
the complement, coindex the subject with the gap,
and (via \ft{stop-gap}) empty \ft{gap} on the mother
\end{itemize}

\myslide{More Filler-Gap Phenomena}

\begin{itemize}
\item Sentences with subject gaps
\item Gaps in coordinate constructions
\end{itemize}

\section{Subject Gaps}

\myslide{Subject Gaps}
\begin{itemize}
\item The ARP revision only allowed missing
complements.
\item But gaps occur in subject position, too:
  \begin{exe}
    \ex \eng{\ul{This problem}, everyone thought  \blank was too easy.}
  \end{exe}
\item We handle these via a lexical rule that, in effect,
moves the contents of the \ft{spr} list into the \ft{gap} list.
\end{itemize}

\myslide{The Subject Extraction Lexical Rule}

\scalebox{0.9}{\avm{[\type*{pi-rule}
input & < X, [ syn & [ head & [\type*{verb} form & fin ] \\
                          val & [ spr & < Z > ] ] \\
                 arg-st & \tag{A}         ]   >\\
output & <Y,  [ syn & [val & [ spr & < ~~  > ] \\ 
                          gap & < \1 > ] \\
                 arg-st & \tag{A} < \1, \ldots >        ]   > ]
}}

\begin{itemize}
\item NB: This says nothing about the phonology, because the
default for \val{pi-rule}s is to leave the phonology unchanged.
\end{itemize}

\myslide{A Lexical Sequence This Licenses}

\avm{< \textnormal{likes}, \ [\type*{word}
      syn & [
      head & [\type*{verb} form & fin ] \\
      % [\type*{agr-cat} agr & [ per & 3 \\ num & sg ] ]
      val & [ spr & <  >\\
              comps & < \2  > ] \\
      gap & < \1 NP [ case & nom \\ agr & 3sing ] > \\
      stop-gap & < ~~ >]\\
      arg-st &      < \1, \2 NP:acc  >   
      ]>
}
    

\begin{itemize}
\item Note that the ARP is satisfied
\end{itemize}


\myslide{A Tree with a Subject Gap}
\MyLogo{\eng{Kim$_i$ we know  \blank$_i$ likes Bobby}}
\scalebox{0.8}{\begin{tikzpicture}
\tikzset{level distance=8ex}
\tikzset{frontier/.style={distance from root=40ex}}
\tikzset{every tree node/.style={align=center,anchor=north}}
\Tree [.\avn{S}{gap & < >} 
         [.\avn{NP}{gap & < >} Kim ]
         \edge[style=double,very thick];
         [.\avn{S/NP}{gap & < NP >} 
           [.\avn{NP}{gap & < >} we ]
           \edge[style=double,very thick];
           [.\avn{VP/NP}{gap & < NP >} 
           \edge[style=double,very thick];
              [.\avn{V}{gap & < >} know ]
              [.\avn{S/NP}{gap & < NP >} 
              \edge[style=double,very thick];
                [.\avn{V/NP}{gap & < NP >} likes ]
                [.\avn{NP}{gap & < >} Bobby ]  ] ] ] ]
\end{tikzpicture}}
                    

\qaslide{Which Node Has a Non-Empty \ft{stop-gap}?}{
\MyLogo{\eng{Kim$_i$ we know  \blank$_i$ likes Bobby}}
\begin{itemize}
\item In the tree on the previous slide only one node has a
  non-empty \ft{stop-gap}.  Which, and why?
  \answer{\\ {\small\emp{The S \eng{we know \blank\ likes Bobby}, the
    head daughter of the Head-Filler Rule: the rule gives it
    [\ft{stop-gap} $\langle$ NP $\rangle$], the same NP as its
    \ft{gap}, so the mother's \ft{gap} is empty.  Every other head
    daughter is a word with the default empty \ft{stop-gap} or is
    constrained to \el\ by the Head-Specifier Rule.}}}
\item And the gapped \lex{likes}: which rule put the NP in its
  \ft{gap}?
  \answer{\\ {\small\emp{The Subject Extraction Lexical Rule, not the ARP:
    the ARP only moves complements into \ft{gap}.}}}
\end{itemize}
}

\section{Coordinate Structure Constraint}

\myslide{Island Constraints}
\begin{itemize}
\item There are configurations that block filler-gap
dependencies, sometimes called \txx{islands}
\item Trying to explain them has been a central topic of
syntactic research since the mid-1960s
\item We'll look at just one, Ross's so-called
\txx{Coordinate Structure Constraint}
\item Loose statement of the constraint: a constituent
outside a coordinate structure cannot be the filler
for a gap inside the coordinate structure
\end{itemize}

\myslide{Coordinate Structure Constraint Examples}

\begin{exe}
  \ex \bad{This problem, nobody finished the extra credit and \blank.}
  \ex \bad{This problem, nobody finished \blank and the extra credit.}
  \ex \bad{This problem, nobody finished \blank and started the extra credit.}
  \ex \bad{This problem, nobody started the extra credit and finished \blank.}
  \ex \eng{This problem, everybody started \blank and nobody finished \blank.}
\end{exe}

\begin{itemize}
\item In a coordinate structure,
  \begin{itemize}
  \item No conjunct can be a gap (\txx{conjunct constraint})
  \item No gap can be contained in a conjunct if its filler is
    outside of that conjunct (\txx{element constraint})
  \item[\ldots] unless each conjunct has a gap that is paired
with the same filler (\txx{across-the-board exception})
\end{itemize}
\end{itemize}

\myslide{These Observations Cry Out for Explanation}
\begin{itemize}
\item In our analysis, the conjunct constraint is an immediate
consequence: individual conjuncts are not on the \ft{arg-st} list
of any word (\lex{and} has an empty \ft{arg-st}), so they can't be
put on the \ft{gap} list.
\item The element constraint and ATB exception suggest that \ft{gap}
is one of those features (along with \ft{val} and \ft{form}) that
must agree across conjuncts.
\item Note: There is no ATB exception to the conjunct constraint.
  \begin{exe}
    \ex \bad{This problem, you can compare only \blank and \blank.}
  \end{exe}
\end{itemize}

\myslide{Our Coordination Rule, So Far}

\scalebox{0.75}{\avm{[ form & \1 \\ val & \0 \\ ind & $s_0$ ]
    \ \into \ 
  [ form & \1 \\ val & \0 \\ ind & $s_1$  ] \ \ldots \  
   [ form & \1 \\ val & \0  \\ ind & $s_{n-1}$ ]
    [ head & conj \\
     ind  & $s_0$ \\ 
    restr & < [ args & < $s_1, \ldots, s_{n-1}, s_n$ > ] >  ]\ \ 
  [ form & \1 \\ val & \0 \\ ind & $s_n$ ]
}}


\begin{itemize}
\item Recall that we have tinkered with what must agree across
conjuncts at various times.
\item Now we'll add \ft{gap} to the things that conjuncts must share.
\end{itemize}


\qaslide{Our Final Coordination Rule}{

\scalebox{0.75}{\avm{[ form & \1 \\ val & \0 \\ ind & $s_0$ \\ gap & \tag{A} ]
    \ \into \ 
  [ form & \1 \\ val & \0 \\ ind & $s_1$  \\ gap & \tag{A}] \ \ldots \  
   [ form & \1 \\ val & \0  \\ ind & $s_{n-1}$ \\ gap & \tag{A}]
    [ head & conj \\
     ind  & $s_0$ \\ 
    restr & < [ args & < $s_1, \ldots, s_{n-1}, s_n$ > ] >  ]\ \ 
  [ form & \1 \\ val & \0 \\ ind & $s_n$ \\ gap & \tag{A}]
}}

\begin{itemize}
\item We've just added \ft{gap} to all the conjuncts and the mother.
\item This makes the conjuncts all have the same gap (if any).
\item Why do we need it on the mother?
  \answer{\\ {\small\emp{Coordination is not headed, so the \ft{gap} Principle doesn't pass the gap up; without it the gap would vanish, and \bad{Kim bought and Sandy sold} would be a sentence.}}}
\end{itemize}
}

\myslide{Closing Remarks on LDDs}
\begin{itemize}
\item This is a huge topic; we've only scratched the
surface
\item There are many more kinds of LDDs, which
would require additional grammar rules
\item There are also more island constraints, which also
need to be explained
\item Our account of the coordinate structure constraint
(based on ideas of Gazdar) is a step in the right
direction, but it would be even better to explain why certain
features must agree across conjuncts
\end{itemize}


\myslide{Overview of LDD}
\begin{itemize}
\item Long-distance dependencies: some examples, and what is new and
  different about them
\item Our approach: the feature \ft{gap}
\item Filling gaps: \ft{stop-gap}, the Head-Filler Rule and
  \lex{easy}-adjectives
\item Subject gaps
\item Coordinate Structure Constraint
\item Problem Set 11: SWB 14.1, 14.2, 14.3
\end{itemize}

\section[Problem Set 11: SWB 14.1, 14.2, 14.3]{Problem Set 11}

\myslide{P0: Semantics Are \lex{easy}}

Add the semantics to the lexeme \lex{easy} given on slide \pageref{l:easy}.

Then give the full \ft{restr} list for the top node (i.e. the whole
sentence) for (\ref{s:easyi}) and (\ref{s:easyii}).  What is the deep
subject of \lex{easy} in each sentence?  

 \begin{exe}
    \ex \label{s:easyi}\eng{My house is easy to find.}
    \ex \label{s:easyii}\eng{Pat is easy to talk to.}
  \end{exe}

\myslide{P1: A Tree with a Gap}
\MyLogo{Based on \citet[Chapter 14, Problem 1]{sag:wasow:bender:2003}}
Draw a tree for (\ref{s:toy-exxb}).  
%You should provide about the same
%level of formal detail as in (\ref{gapbb-tree}).
Use abbreviations for the node labels, and show the value of \ft{gap}
on all nodes.  Show the value of \ft{stop-gap} on any node where it is
non-empty.

\begin{exe}
  \ex \label{s:toy-exxb} \eng{This baby, I know that they handed a toy to {\blank}.}
\end{exe}

\myslide{P2: Blocking Filled Gaps}
\MyLogo{Based on \citet[Chapter 14, Problem 2]{sag:wasow:bender:2003}}

Examples (i) and (ii) are well-formed, but example (iii) is
ungrammatical:

\begin{exe}
\exi{(i)}  \eng{Pat thinks that I rely on some sort of trick.}
\exi{(ii)} \eng{This mnemonic, Pat thinks that I rely on.}
\exi{(iii)} \bad{This mnemonic, Pat thinks that I rely on some sort of trick.}
\end{exe}
Explain in detail why the mechanisms that license (i) and (ii) do not
also permit (iii).

\myslide{P3: Subject Gaps} 
\MyLogo{Based on \citet[Chapter 14, Problem 3]{sag:wasow:bender:2003}} 

This problem is to make sure you understand how our
analysis accounts for examples like (\ref{se-s}).

\begin{exe}
\ex \label{se-s}
\begin{xlisti}
\ex \label{se-sa}\eng{Which candidates do you think like %raw oysters?
oysters on the half-shell?}
\ex \label{se-sb} \eng{That candidate, I think likes %raw oysters.
oysters on the half-shell.}
\end{xlisti}
\end{exe}

\begin{itemize}
\item [A.] Sketch the family of lexical sequences for \lex{likes} that is the input
to the Subject Extraction Lexical Rule.

\item [B.] Sketch the family of lexical sequences for \lex{likes} that is the
corresponding output of the Subject Extraction Lexical Rule.

%\item [(b)] Sketch the lexical entry for {\it likes} that is the output
%of the Subject Extraction Lexical Rule.

\item[C.] Sketch the %entire 
tree for the sentence in (\ref{se-sb}).
%Provide about the same level of formal detail as in (\ref{gapbb-tree}).
Use abbreviations for node labels, but show the value of \ft{gap} on
all nodes and the value of \ft{stop-gap} on any node where it is non-empty.
You may abbreviate the structure over the NP \eng{oysters on the
half-shell} with a triangle.
%Be sure to explain how our grammar ensures that all relevant feature
%specifications are as they are in your sketch.

\item[D.] Does our analysis correctly predict the contrast between
(\ref{se-sb}) and (\ref{onew})?\\
\begin{exe}
  \ex \label{onew} \bad{Those candidates, I think likes %raw oysters.
    oysters on the half-shell.}
\end{exe}
Explain why or why not.
\end{itemize}

% \myslide{Questions (2017)}
% \input{lec-11-ldd-q}

\myslide{Acknowledgments and References}

\begin{itemize}
\item Course design and slides borrow heavily from Emily Bender's course:
\textit{Linguistics 566: Introduction to Syntax for Computational Linguistics}
\\ \url{http://courses.washington.edu/ling566}
\item This lecture complements \citet[Chapter 14]{sag:wasow:bender:2003}
\end{itemize}


\bibliographystyle{aclnat}
\bibliography{abb,mtg,nlp,ling}

\end{document}
%%% Local Variables: 
%%% coding: utf-8
%%% mode: latex
%%% TeX-PDF-mode: t
%%% TeX-engine: xetex
%%% End: 

